\documentclass{uflamon}          % classe base para a monografia

%==============================================================================
% Utilizacao de pacotes
\usepackage[T1]{fontenc}         % usa fontes postscript com acentos
\usepackage[brazil]{babel}       % hifenização e títulos em português do Brasil
\usepackage[utf8]{inputenc}     % permite edição direta com acentos
\usepackage{amsmath}             % pacote da AMS para Matemática Avançada
\usepackage{amssymb}             % símbolos extras da AMS
\usepackage{latexsym}            % símbolos extras do LaTeX
\usepackage{graphicx}            % para inserção de gráficos
\usepackage{listings}            % para inserção de código
\usepackage{fancyvrb}            % para inserção de saídas de comandos
%\usepackage{enumerate}           % para personalizar lista enumeradas 
											%(incluso na classe)
\usepackage{longtable}           % para tambelas muito grandes NOVO!!!!

\usepackage{colortbl} % cores em tabelas
\newcolumntype{Z}{|>{\columncolor[gray]{0.9}}l|} %cor cinza em células
%\usepackage{array} % já incluso na classe
\newcolumntype{L}[1]{>{\raggedright\let\newline\\\arraybackslash\hspace{0pt}}m{#1}}
\newcolumntype{C}[1]{>{\centering\let\newline\\\arraybackslash\hspace{0pt}}m{#1}}
\newcolumntype{R}[1]{>{\raggedleft\let\newline\\\arraybackslash\hspace{0pt}}m{#1}}
\usepackage{multirow} % para juntar duas linhas em uma só

\usepackage{multicol} % para uso de várias colunas

% cores para os links cruzados
\usepackage{color}
\definecolor{rltred}{rgb}{0.2,0,0}
\definecolor{rltgreen}{rgb}{0,0.2,0}
\definecolor{rltblue}{rgb}{0,0,0.2}

\usepackage[colorlinks=true,
            urlcolor=rltblue,       % \href{...}{...} external (URL)
            filecolor=rltgreen,     % \href{...} local file
            linkcolor=rltred,       % \ref{...} and \pageref{...}
            citecolor=rltgreen,
            pdftitle={Exemplo de Uso da Classe Uflamon},
          pdfauthor={Joaquim Quinteiro Uchôa},
          pdfsubject={Este texto tem por objetivo servir de exemplo da classe Uflamon.},
          pdfkeywords={Comunicação Científica. 2. Pesquisa . 3. Pesquisa Científica. 
 					 4. Redação. 5. Monografia.}%
]{hyperref} % para referência cruzadas
%\usepackage{hyperref}            % para referência cruzadas
\usepackage{subfigure}           % figuras dentro de figuras
\usepackage{caption}            % remodelando o formato dos títulos de 
                                 % tabelas e figuras

% configuração padrão do listings   
\lstset{
   language=Java,
   extendedchars=true,
   tabsize=3,
   basicstyle=\footnotesize\ttfamily,
   stringstyle=\em,
   showstringspaces=false 
}

% para referências de acordo com a ABNT
% precisa instalar o abntex2 antes!!!
% http://abntex.codigolivre.org.br/
% comente se pretende usar outro padrão

%abnt-emphasize=bf coloca o título das bibliografias em negrito
%abnt-thesis-year=both
\usepackage[alf,abnt-etal-cite=3,abnt-etal-list=3,abnt-url-package=url,abnt-emphasize=bf]{abntex2cite}

% evite usar o hyperref com abntex, pode dar caca em urls... no linha anterior, informo
% para incluir urls usando o pacote url e não o hyperref
%
% caso queira o hyperref com abntex, comente a linha anterior e descomente a seguinte
%\usepackage[alf,abnt-etal-cite=3,abnt-etal-list=0,abnt-etal-text=emph]{abntex2cite}
%
% caso vc ainda use a versão anterior da abntex, comente a linha incluindo o abntex2cite
% e descomente a próxima linha 
%\usepackage[alf,abnt-etal-cite=3,abnt-etal-list=0,abnt-etal-text=emph]{abntcite}


% redefinindo formatação de títulos de tabelas e figuras


%==============================================================================
% para os fãs do Word, descomente as linhas abaixo
%\sloppy %mais espaço entre as linhas
%\usepackage{identfirst} %identando-se a primeira linha de cada seção
%\noindentfirst % Tire o comentário para manter o padrão do LaTeX.

%==============================================================================
% definido comandos na monografia - não é necessário na sua monografia 
% apenas para exemplificar a definição de novos comandos
\newcommand{\defs}[1]{\textsl{#1}}


% Especificando hifenizações que por ventura LaTeX não saiba fazer
% Por padrão 99,9% dos termos em português devem ser hifenizados corretamente.
\hyphenation{hardware software Li-nux am-bien-te diag-nos-ti-car coor-de-na-ção 
FAE-PE Recovery TelEduc Williams UFLA}

%==============================================================================
% Dados da monografia, capa: autor, titulo, banca, etc... - SUBSTITUA DE ACORDO
%==============================================================================
\author{Joaquim Quinteiro Uchôa}
\title{Uso da Classe Uflamon}
\subtitle{Exemplo para os Usuários}
\engtitle{Use of Uflamon Class}
\engsubtitle{Sample for Users}
\edicao{3$^a$ edição revista, atualizada e ampliada}
\date{2016}
\tipo{Tese apresentada à Universidade Federal de Lavras, como parte das exigências do Programa de Pós-Graduação em Monografia, área de concentração em TCC, para a obtenção do título de Doutor.}
% use \orientador ou \orientadora quando for o caso
\orientador{Prof. DSc. José Orientador}
%\orientadora{}
% use \coorientador ou \coorientadora quando for o caso
\coorientadora{Prof. DSc. Maria Orientadora } % comente se não tiver coorientador
%\coorientador{}
\local{Lavras -- MG}
\bancaum{Prof. MSc. Antônio Banca Um}{UFM}
\bancadois{Prof. DSc. João Banca Dois}{FCO} % comente se sua banca tiver só um professor
\bancatres{Profa. Esp. Eliza Banca Três}{BELMIS}
\bancaquatro{Prof. Esp. Carlos Banca Quatro}{IBGPLUS}
\defesa{30 de Fevereiro de 2016}
%==============================================================================
%##################################################
% Dados para Ficha catalográfica, gerada pelo sistema da Biblioteca da UFLA
% http://www.biblioteca.ufla.br/FichaCatalografica/
% dados para ficha catalográfica
% Elaboração da Ficha Catalográfica
\preparofichacat{Ficha catalográfica elaborada pela Coordenadoria de Processos Técnicos \\ da Biblioteca Universitária da UFLA}
% primeiro autor - como na primeira linha da ficha catalográfica
\fcautor{Uchôa, Joaquim Quinteiro}
% autores, separados por vírgula - na ficha catalográfica, no formato que
% vem após o título e a barra ("/")
\fcautores{Joaquim Quinteiro Uchôa}
% caso trabalho seja ilustrado (figuras, gráficos, tabelas, etc.), 
% então informar por meio do comando a seguir
% caso não seja ilustrado, basta comentá-lo
\fcilustrado{il.}
% dados da edição para a ficha 
\fcedicao{2$^a$ ed. rev., atual. e ampl.}
% tipo do trabalho (tese, dissertação, etc.), de acordo com sistema
% de geração de ficha catalográfica
\fctipo{Tese(doutorado)}
% ano da defesa, só precisa informar se for diferente do ano da publicação
% se forem iguais, comente a linha a seguir
\fcdatadefesa{2016}
% preencher aqui com os dados de catalogação gerados pelo sistema
\fccatalogacao{1. TCC. 2. Monografia. 3. Dissertação. 4. Tese. 5. Trabalho Científico – Normas. I. Universidade Federal de Lavras. II. Título.}
\fcclasi{808.066}

%##################################################

%\antesfichacat{\noindent Para citar este documento: \\UNIVERSIDADE FEDERAL DE LAVRAS. Biblioteca Universitária. \textbf{Manual de normalização e estrutura de trabalhos acadêmicos: TCC, monografias, dissertações e teses}. 2. ed. rev., atual. e ampl. Lavras, 2015. Disponível em: \url{http://www.biblioteca.ufla.br/wordpress/wpcontent/uploads/bdtd/manual_normalizacao_UFLA.pdf}. Acesso em: data de acesso.}

%\depoisfichacat{\noindent A reprodução e a divulgação total ou parcial deste trabalho são autorizadas, por qualquer meio convencional ou eletrônico, para fins de estudo e pesquisa, desde que citada a fonte.\\
%\newline
%{\small Este documento possui páginas em branco para facilitar a impressão frente-e-verso.}}

%##################################################

%##################################################

% para os exemplos do manual
%\newenvironment{exemplomanual}{
%\vspace{0.5cm}
%\noindent\begin{minipage}{\textwidth}
%\noindent\rule{\textwidth}{0.5pt}
%\vspace{-1cm}
%\begin{flushleft}
%}{
%\end{flushleft}
%\vspace{-0.6cm}
%\noindent\rule{\textwidth}{0.5pt}
%\vspace{0.3cm}
%\end{minipage}
%}

%\newenvironment{exemplomanuallista}{
%\vspace{0.3cm}
%\noindent\begin{minipage}{\textwidth - 0.5cm}
%\noindent\rule{\textwidth}{0.5pt}
%\vspace{-1cm}
%\begin{flushleft}
%}{
%\end{flushleft}
%\vspace{-0.6cm}
%\noindent\rule{\textwidth}{0.5pt}
%\vspace{0.3cm}
%\end{minipage}
%}

% por conta de alguns exemplos
%\usepackage{setspace}

%##################################################

% se vc já defendeu e tem o arquivo escaneado da folha de rosto, 
% descomente e altere o nome do arquivo
%\folhaAprovacaoAssinada{folharosto}

% Aqui começa o documento propriamente dito
\begin{document}

\maketitle

\dedic{Dedico aos meus pais, Elisa e Gilberto, ao meu irmão Bruno,e a todos aqueles que me inspiraram durante essa caminhada.}     % Dedicatórias


\thanks{
Essa era a página que eu tanto desejava escrever. Uma  representação de mais uma etapa que está sendo concluída, de um novo ciclo que está prestes a começar. Irei chamar essa de página de página da gratidão. Há muitas pessoas que sou grato por ter chegado até aqui!
Meus pais, Gilberto e Elisa, por abrirem mão de muitas coisas, desde o início, para investir na minha educação, nos meus estudos. Ainda criança era bem difícil de entender o porquê, mas hoje tenho absoluta certeza de que cada palavra de incentivo dos dois foi o que me trouxe até aqui! Meu irmão, Bruno, por estar sempre presente, contagiando ainda mais os meus dias de alegria!
À minha família, obrigado pelo carinho! Aos meus avós, pelas rezas, pelo cuidado e preocupação sobre meu bem estar. O neto está quase voltado!

Durante minha ida à Lavras, muitos amigos ficaram. Obrigado pela paciência e por entender essa distância. Luiza Fujii e Marcell: estão sempre em meus pensamentos! Dina, Carol, Arthur e Gabriel: obrigado por todo o apoio durante esse último ano!

Durante minha trajetória, tive o prazer de fazer muitos amigos. Minhas amigas da República Mesadose: agradeço pelas festas! Aos meus amigos do Palácio da Cartas, obrigado pela companhia!\

Ao escolher Ciência da Computação, meu objetivo durante o curso foi, além de aperfeiçoar em assuntos de meu interesse, foi de, alguma forma , transmitir os conhecimentos que adquiri durante toda a jornada. Sou muito grato por todos os grupos de estudos, pesquisa e extensão que passei, entre eles, CACC, Comp Júnior, PETI, Alcance, DGTI, SETI. Extende-se a todos os amigos que fiz em cada um deles. Agradeço ao Google Developers Group, pela confiança em poder tocar esse projeto em Lavras, tentando tornar a comunidade de desenvolvedores da cidade mais forte e por ajudar a incentivar o empoderamento feminino na tecnologia.
 
Uma boa parte da minha graduação foi marcada pelo intercâmbtio nos EUA, e lá tive o imenso prazer de conhecer pessoas que foram fundamentais e que cada um tem todo o meu respeito. Agradeço aos professores da Auburn University at Montgomery e da Arizona State Univesity. E, principalmente, os amigos que conquistei por lá. Foram aprendizados para a vida toda. Todo dia era dia de dizer "It's gonna be the best day of my life", e foi.

Agradeço aos mestres e doutores que foram meus professores na Universidade Federal de Lavras. Todo o conhecimento que adquiri foram essenciais por eu ser quem eu sou. Obrigado pela dedicação e paciência!

Em especial, agradeço o professor André Pimenta Freire, por me orientar em diversos projetos na universidade. Encontrei durante os projetos o que realmente eu gosto de fazer e estudar, e o que quero dedicar boa parte dos meus estudos daqui em diante: a experiência do usuário.

Agradeço aos meus companheiros de trabalho da Progolden, onde realizei o meu estágio, resultando o trabalho descrito neste documento. Aprendi a rabalhar em grupo vi de perto como um projeto sai do papel e se transforma em um software. E e assim, daqui em diante, que espero que a vida me msorte vários projetos sendo executados e que, principalmente, mude a vida das pessoas.

Assim, dedico à todos(as) os(as) corajosos(as) que colocaram na cabeça a ideia de fazer Ciência da Computação!
}        


	
	
	% Agradecimentos

\epigrafe{ % citação opcional
What you want, what you need.\\
And what'd you come here for?\\
Raise up to your ability\\
You gotta focus on your ability!\\
(Houdini - Foster the People)}

% palavras-chave
\palchaves{Resumo. Palavras. Representativas.}
\resumo{O resumo deve conter palavras representativas do conteúdo do trabalho, localizadas abaixo do resumo, separadas por dois espaços, antecedidas da expressão palavras-chave. Essas palavras representativas são grafadas com a letra inicial em maiúscula, separadas entre si por ponto.}  % Resumo (digite aqui o resumo)

% keywords devem vir antes do abstract
\keywords{Summary. Words. Representative.} % keywords
\abstract{The abstract should contain representative words of the work content, located below the abstract, separated by two spaces, preceded by the keyword expression. These representative words are spelled with the first letter capitalized, separated by point.}

%##################################################

% Dados do guia
%\begin{titlepage}
%\pagestyle{empty}
%\renewcommand{\baselinestretch}{1}
%\enlargethispage{1.5cm}
%\input{reitoria}
%\cleardoublepage
%\end{titlepage}

%##################################################

% descomente para habilitar a lista desejada
\listoffigures                             % Lista de Figuras
%\listofilustracoes
%\listofgraficos							   % Lista de Gráficos
\listoftables                              % Lista de Tabelas
\listofquadros							   % Lista de Quadros
%\listofexemplos
%\listofteoremas
\tableofcontents                           % Sumário

\clearpage

\pagestyle{ufla}

%==============================================================================
% incluindo os capitulos
\chapter{INTRODUÇÃO}

%TODO Inclua na introdução também do que se trata a empresa e tipo de projeto que você trabalhou (Done)
O objetivo deste documento é apresentar as tarefas realizadas durante o estágio na empresa Progolden Soluções Tecnológicas LTDA, uma startup fundada em 2013 com foco no desenvolvimento de soluções tecnológicas inovadoras. O presente documento descreve conhecimentos obtidos nos estudos de Design Centrado no Usuário, Usabilidade, Experiência do Usuário e \textit{Lean UX} (\textit{User Experience}) durante o estágio, aplicados em um projeto de desenvolvimento de software para planejamento estratégico.

O foco do Design Centrado no Usuário (DCU) é estudar quem são os usuários, quais as suas necessidades e como desenvolver soluções para eles. Segundo a ISO 9241-210 \cite{iso9241}, que descreve processos de projeto centrados no ser humano para sistemas interativos; design centrado no usuário é uma atividade multidisciplinar, que incorpora fatores humanos e conhecimentos ergonômicos e técnicas com o objetivo de aumentar a eficácia e produtividade, melhorar as condições de trabalho humano e neutralizar os possíveis efeitos adversos da utilização na saúde humana, segurança e desempenho. O desenvolvimento do produto, então, teve foco no usuário e como este poderá se beneficiar da melhor forma possível. 

%TODO  Está com transcrição literal, tal qual escrito na ISO? se sim, precisa vir entre aspas

Ainda segundo a ISO 9241-210 \cite{iso9241}, existem quatro atividades de projeto centradas no usuário que precisam ser iniciadas nas primeiras etapas de um projeto, que são:

\begin{itemize}

\item "entender e especificar o contexto de uso";

\item "especificar os requisitos do usuário e da organização";

\item "produzir soluções de design";

\item "avaliar e homologar".

\end{itemize}

Juntamente com o DCU, a metodologia de trabalho do projeto foi o Desenvolvimento Ágil. Atualmente, as equipes enfrentam uma intensa pressão de concorrentes que estão usando técnicas com o desenvolvimento de software ágil, integração contínua e implantação contínua para reduzir radicalmente seus tempos de ciclo \cite{gothelf2013lean}. O manifesto ágil para desenvolvimento de software valoriza os indivíduos, interações, software em funcionamento, colaboração com o cliente, e respostas às mudanças \cite{robbinsd:beck2001agile}. 

Durante o estágio, percebeu-se que, intuitivamente, juntamente com a metodologia ágil aplicada, estava-se trabalhando com a metodologia Lean UX (User Experience). 

O \textit{Lean UX} é definido como uma abordagem para um desenvolvimento de software centrado no usuário, especialmente em startups, criando produtos radicalmente novos. São identificadas três principais influências no \textit{Lean UX}: movimento design thinking, método Learn startup e desenvolvimento de software ágil  \cite{gothelf2013lean}.

O  estágio teve como objetivo aplicar conceitos de design centrado no usuário e usabilidade em um software de planejamento estratégico, capacitando o estagiário para o desenvolvimento de soluções, relizando funções de analista de sistemas, desenvolvedor \textit{front-end}, testador e \textit{user experience researcher}. As principais atividades atribuídas durante o estágio foram:

\begin{itemize}

\item realização de análise de requisitos;

\item especificação de sistemas;

\item desenvolvimento de protótipos;

\item realização de avaliação heurística;

\item realização de testes e correções de erros.

\end{itemize}

O restante desse documento está organizado da seguinte forma: o Capítulo 2 apresenta a descrição do local de realização do estágio.  No Capítulo 3, são apresentadas as técnicas e processos utilizados.  O Capítulo 4 apresenta as atividades desenvolvidas e, finalmente, o Capítulo 5 apresenta as conclusões.

\chapter{ELEMENTOS DO TEXTO}
\label{cap:elementos}

Este capítulo apresenta o uso básico de equações, figuras e tabelas no código da monografia, bem como o posicionamento das legendas, segundo as normas da UFLA.

\section{Utilizando Recursos do \LaTeX}

\subsection{Inserindo Comandos Definidos}

Esta subseção apresenta o uso de comandos definidos pelo usuário no preâmbulo do arquivo principal \LaTeX e alguns exemplos do modo matemático. Por exemplo, na texto a seguir é utilizado o comando  \verb+\defs+, definido anteriormente pelo próprio autor do texto:

\begin{quote}
Os conjuntos fundamentais da teoria são os \defs{conjuntos elementares}. Se $E$ é um conjunto elementar, $des(E)$ denota a descrição dessa classe de equivalência. Essa descrição é função do conjunto de atributos que define $R$. Note que, dados $x,y \in E$, onde $E$ é um conjunto elementar em $A$, $x$ e $y$ são indiscerníveis, i.e., no espaço $A=(U,R)$ não se consegue distinguir $x$ de $y$, pois $des(x) = des(y) = des(E)$. 
\end{quote}

\subsection{Inserindo Figuras}

A Figura~\ref{fig:exemplo} é apenas um exemplo de figura para que o usuário da classe possa saber como uma figura pode ser inserida e referenciada automaticamente no texto.É importante observar que legendas de figuras ficam abaixo de seu conteúdo.

\begin{figure}[!htb]
\centering
\caption{Uma Figura de Exemplo} %legenda
\includegraphics[scale=0.9]{gradpenguin}\\  % o 0.9 indica 90% do tamanho original
% pdfLaTeX aceita figuras no formato PNG, JPG ou PDF
% figuras vetoriais podem ser exportadas para eps e depois convertidas para pdf usando epstopdf
{\small Fonte: fonte da figura} %Fonte da imagem
\label{fig:exemplo} %rotulo para refencia
\end{figure}

\subsection{Inserindo Saídas de Comandos e Código}

A menos que sejam trechos pequenos, saídas de comandos, trechos de arquivos de configuração e código de aplicativos devem ser inseridos como figura, como mostrado, respectivamente, na Figura~\ref{fig:exemplocomando}, Figura~\ref{fig:exemploconfig} e Figura~\ref{fig:exemplocodigo1}. Para comandos e configuração, recomenda-se o uso do pacote {\tt fancyvrb}, o que pode ser visto na Figura~\ref{fig:exemplocomando} e Figura~\ref{fig:exemploconfig}.

Para inserção de código, recomenda-se o uso do pacote {\tt listings}, que permite melhor apresentação do mesmo, o que é mostrado na Figura~\ref{fig:exemplocodigo1}. Além disso, esses dois pacotes permitem a inserção de texto/código em arquivos externos, sem inclusão direta no arquivo \LaTeX. Isso pode ser verificado no exemplo de uso do {\tt listings} apresentado na Figura~\ref{fig:exemplocodigo2}

\begin{figure}[!htb]
\centering
\caption{Inserindo Comando} %legenda
\begin{Verbatim}[fontsize=\small]
$ dir monografia*
-rw-r--r--  1 joukim users   3650 Set 12 17:56 monografia.aux
-rw-r--r--  1 joukim users   6366 Set 12 17:43 monografia.bbl
-rw-r--r--  1 joukim users    290 Set 12 17:56 monografia.lof
-rw-r--r--  1 joukim users  27937 Set 12 17:56 monografia.log
-rw-r--r--  1 joukim users    194 Set 12 17:56 monografia.lot
-rw-r--r--  1 joukim users    715 Set 12 17:56 monografia.out
-rw-r--r--  1 joukim users 159243 Set 12 17:56 monografia.pdf
-rw-r--r--  1 joukim users   4559 Set 12 17:47 monografia.tex
-rw-r--r--  1 joukim users    964 Set 12 17:56 monografia.toc
\end{Verbatim} 
%$ - esse comentário é para não confundir editor de texto
{\small Fonte: fonte da figura} %Fonte da imagem
\label{fig:exemplocomando} %rotulo para refencia
\end{figure}

\begin{figure}[!htb]
\centering
\caption{Inserindo Trecho de Arquivo de Configuração} %legenda
\begin{Verbatim}[fontsize=\small]
// named.conf for Red Hat caching-nameserver
options {
        directory "/var/named";
        dump-file "/var/named/data/cache_dump.db";
        statistics-file "/var/named/data/named_stats.txt";
        // query-source address * port 53;
};
\end{Verbatim} 
{\small Fonte: fonte da figura} %Fonte da imagem
\label{fig:exemploconfig} %rotulo para refencia
\end{figure}


\begin{figure}[!htb]
\centering
\caption{Inserindo Código Diretamente no Arquivo \LaTeX} %legenda
\begin{lstlisting}
// exit the program
public void on_buttonExit_clicked() {
	System.exit(0);
}

// copy data
public void on_buttonCopy_clicked() {
	labelShow.setText(entryData.getText());
}

// print version of Java
public static void main(String[] args) {
	System.out.println(System.getProperty("java.fullversion"));
}
\end{lstlisting} 
{\small Fonte: fonte da figura} %Fonte da imagem
\label{fig:exemplocodigo1} %rotulo para refencia
\end{figure}


\begin{figure}[H]
\centering
\caption{Inserindo Código a Partir do Código-Fonte} %legenda
\label{fig:exemplocodigo2} %rotulo para refencia
\lstinputlisting{Hello.java}
{\small Fonte: fonte da figura} %Fonte da imagem
\end{figure}

\subsection{Inserindo Quadros e Tabelas}

Escrever um quadro ou tabela e referenciá-los é bem simples. Por exemplo o Quadro~\ref{tab:exemplo} ilustra a criação de um quadro, tendo aqui seu referenciamento no texto. É importante observar o posicionamento da legenda antes do corpo da tabela e da fonte após. Outros exemplos são mostrados na Tabela~\ref{tab:outro} e Tabela~\ref{tab:maisum}.

\begin{quadro}[htb]
  \begin{center}
    \caption{Exemplo de Quadro} 
    \label{tab:exemplo}
    \vspace{0.2cm}
    \footnotesize
    \begin{tabular}{|c|c|c|c|c|c|}
      \hline
      $U$ & $vitA$ & $vitC$ & $vitD$ & $prot$ & $lip$ \\
      \hline
      \hline
      $d_1$ & 1 & 3 & 4 & 2 & 3\\
      $d_2$ & 1 & 3 & 3 & 3 & 2\\
      $d_3$ & 1 & 3 & 4 & 3 & 1\\
      $d_4$ & 3 & 5 & 2 & 5 & 2\\
      $d_5$ & 4 & 5 & 2 & 5 & 1\\
      $d_6$ & 3 & 5 & 2 & 3 & 4\\
      $d_7$ & 4 & 4 & 1 & 3 & 2\\
      \hline 
    \end{tabular}
  \end{center}
  \centering {\small Fonte: fonte do quadro} %Fonte do quadro
\end{quadro}

\begin{table}[!htb]
\begin{center}
  \caption{Recursos do {\ttfamily syslog}}
  \label{tab:outro}
  \small
  \begin{tabular}{l|p{9cm}}
    \hline
    \rowcolor[gray]{.9}
    \bf Recurso & \bf {\em Daemons} Associados (Alguns Exemplos) \\
    \hline
    \hline
    \tt kern & \em kernel  \\
    \tt user & processos dos usuários ({\tt ntpd}) \\
    \tt mail & softwares relacionados com o correio eletrônico ({\tt sendmail})\\
    \tt daemon & {\em daemons} do sistema ({\tt gated}, {\tt inetd}, 
    {\tt named}, {\tt ntpd})\\
    \tt auth &  comandos relacionados à autorização e segurança 
    ({\tt login}, {\tt rlogin}, {\tt su}, {\tt passwd}) \\
    \tt lpr & spool de impressão ({\tt lpd})\\
    \tt news & sistema de notícias da usenet ({\tt nnrpd})\\
    \tt uucp & destinado ao {\tt uucp}\\
    \tt cron & relacionado ao {\em daemon} {\tt cron}\\
    \tt mark &  registros de data/hora gerados a intervalos regulares 
    ({\tt syslogd})\\
    \tt local0-7 & 8 tipos de mensagens locais \newline
    ({\tt tcpd -- local7}, {\tt sudo -- local2}, {\tt popper - local0}) \\
    \tt syslog &  mensagens internas ao {\tt syslog}\\
    \tt authpriv & mensagens privadas de autorização\\
    \tt ftp & associado ao {\tt ftpd} ({\em daemon} do {\tt ftp}) \\
    \tt * &  todos os recursos com exceção do {\tt mark}\\
    \hline
  \end{tabular}
\end{center}
\centering {\small Fonte: fonte da tabela} %Fonte da tabela
\end{table}

\begin{table}[!htb]
\caption{Opções do Comando {\ttfamily at}}
\label{tab:maisum}
\begin{center}
  \small
  \begin{tabular}{l|p{9cm}}
    \hline 
    \rowcolor[gray]{.9}
    \bf Opção & \bf Descrição\\
    \hline
    \hline 
    \tt -c & exibe os jobs registrados\\
    \tt -d & remove um job específico\\
    \tt -f & permite que os comandos sejam lidos a partir de um arquivo (não pela
    entrada-padrão)\\
    \tt -l & lista todos os jobs associados a um usuário\\
    \tt -m & envia um e-mail ao usuário quando o job for finalizado\\
    \hline \end{tabular}\\
\end{center}
\centering {\small Fonte: fonte da tabela} %Fonte da tabela
\end{table}

\subsection{Inserindo Equações}

Equações devem ser numeradas, com a numeração, em parênteses à direita da mesma. Isso é ilustrado na Equação~\ref{eq:exemplo}.

\begin{equation}
\label{eq:exemplo}
f'(x) = \int^{x^2}_{x^{-1}} xdx 
\end{equation}


\section{Usando Referências}
A equipe do curso não impõe normas rígidas para o formato a ser adotado nas referências bibliográficas, desde que seja usado um padrão acadêmico conhecido. Caso os autores não possuam um padrão preferencial, recomenda-se o formato estipulado pela ABNT \cite{NBR6023:2002}. A Biblioteca Central da UFLA disponibiliza um manual \cite{BIBUFLA2001} que orienta o uso dessas normas. Se os autores estiverem utilizando \LaTeX, recomenda-se o uso do pacote Abn\TeX\footnote{Disponível em \url{http://abntex.codigolivre.org.br/}.} \cite{Weber2003}. 

Obviamente, recomenda-se a leitura de textos sobre a escrita acadêmica e produção de trabalhos de conclusão para garantir não só qualidade estética e de formatação, mas também de conteúdo. Entre outros, pode-se recomendar a leitura de \cite{Silva2005}, \cite{Martins2000}, \cite{Gil2002}, \cite{Franca2001}, \cite{Eco1996}, \cite{Moura1998}, \cite{Booth2000}, \cite{Hexsel2004}, \cite{Porto2002} e \cite{Henz2003}. 


\chapter{CONCLUSÃO}

O estágio realizado na Progolden, como trabalho de conclusão de Curso para obtenção do título de Bacharel em Ciência da Computação, proporcionou a aplicação dos conhecimentos adquiridos no período acadêmico, em um contexto de mercado de trabalho e em um projeto de nível nacional.

O estágio proporcionou conhecimentos, tanto em nível acadêmico quanto em nível profissional. Por se tratar de uma \textit{startup} fundada por alunos e professores do Departamento de Ciência da Computação da Universidade Federal de Lavras, o estágio contribuiu para ter uma visão de como todas as áreas da empresa se relacionam e como a empresa surgiu, a partir de ideias desenvolvidas dentro da própria universidade.

%tirar essa parte do plano de desenvolvimento institucional
A participação no projeto de planejamento estratégico possibilitou experiência de desenvolver um produto, no caso, um software, desde sua concepção até a entrega. 

Os conhecimentos adquiridos na disciplina Engenharia de software ajudaram na realização das atividades de levantamento de requisitos, análise de sistemas e documentações do software.

As disciplinas Algoritmos e estruturas de dados e Banco de dados proporcionaram uma base para o desenvolvimento web e a lógica de programação, além de ajudar o estagiário no entendimento completo do software e na correção de problemas encontrados.

Os conhecimentos da disciplina Interface Homem-máquina, como design centrado no usuário, usabilidade e avaliação heurística foram aplicados durante o estágio, possibilitando o estagiário praticar metodologias que serão importantes no mercado de trabalho.

%ta mais pra médio porte kkk
A participação no projeto trouxe conhecimento para a gerência de projetos, como a divisão de atividades entre os membros e prazos de entrega, e a possibilidade de trabalhar em um projeto utilizando metodologias ágeis.

É comum que, durante as disciplinas, o aluno fique atrelado à parte conceitual e teórica, e muitas das vezes, não sabe ou encontra uma alternativa de como utilizá-las na prática. Como por exemplo, a utilização de versionamento de código, uma abordagem empregada pela maioria das empresas de desenvolvimento de software e que, durante o curso, é visto apenas na teoria, em sala de aula. A prática pôde ser exercida durante a realização do estágio.

%ta falando que é pras universidades ?
O estágio mostrou como o curso de Ciência da Computação pode contribuir na resolução de problemas do cotidiano, no caso, a falta de um software de planejamento estratégico que pudesse agilizar o trabalho de entidades organizacionais e empresas. A partir da problemática, a solução foi construída com foco no usuário e em suas reais necessidades.

Ver como a tecnologia pode influenciar, estimular e facilitar os procedimentos sobre outras áreas do conhecimento é uma motivação tamanha para continuar os estudos na área de Ciência da Computação. 

\section{Trabalhos futuros}

Como trabalho futuro, é importante a realização de testes de usabilidade com usuários reais, para que seja possível encontrar os pontos fracos e ajustá-los de acordo com as necessidades, alinhando a entrega com uma usabilidade eficiente. Com o surgimento de novos módulos e funcionalidades do software, é importante realizar novas avaliações heurísticas.

%==============================================================================
% Incluindo bibliografia
%\bibliographystyle{plain}             % estilo para labels em numeros
%\bibliographystyle{alpha}             % estilo para labels em iniciais
\bibliographystyle{abntex2-alf}           % estilo para referências usando ABNT, 
                                       % precisa instalar o abntex para usar!!!

%inclui Referências Bibliográficas
%inclui Referências Bibliográficas
\referencias
\bibliography{refbib}			% arquivo exemplo refbib.bib
%==============================================================================
% Incluindo anexos num1erados com letras maiusculas.
%\apendices
\apendice{Tarefas da Avaliação Heurística}
\label{cap:apendice}

1. Login do Sistema

Nesta tarefa espera-se que o usuário realize o login no sistema.

2.Cadastro de Domínio

Nesta tarefa espera-se que o usuário cadastra o domínio (verifica o quão importante essa tarefa é).

3. Cadastro de estrutura

Nesta tarefa, espera-se que o usuário cadastre uma nova estrutura.

4. Cadastro de novo plano

Nesta tarefa, espera-se que o usuário cadastre um novo plano de metas. 

5.Cadastro de novo plano de metas

Nesta tarefa, espera-se que o usuário cadastre um novo plano de metas, cadastrando todo o fluxo de níveis,  a partir de um plano já cadastrado anteriormente. Será entregue um papel com as informações que devem ser inseridas em cada nível). Realizar o cadastro até a meta (não inserindo valores para meta).


%==============================================================================
% Fim do texto
\end{document}
